\chapter{跨模型与跨病种外部扩展验证}\label{ch:external}

\section{外部验证的设计原则与独立性界定}\label{sec:ext_design}

\subsection{为什么必须开展外部独立场景验证？}
在第五章中，本文基于全美 CDC NHSN 真实的周度流感住院面板数据，对理论仿射模型进行了严格的计量检验。然而，任何严谨的数理统计与经验科学推断都必须直面以下两项方法学质疑：
\begin{enumerate}[label=(\arabic*)]
    \item **对自设基准模型的过拟合风险**：第五章的相对预测损失建立在研究者自设的平滑卷积基准之上。在真实疾控决策实务中，公共卫生机构使用的是集合了数十所顶尖高校与科研机构的复杂前瞻预测模型。理论推导的规模反比递减形态在真实的工业级预测系统下是否依然成立？
    \item **单一病种与单一时间分辨率的局限性**：季节性流感具有相对固定的代际时间（约 3 至 4 天）与周度报告节奏。面对高致病性冠状病毒等不同传播特征的病原体，以及在更高频的日度监测尺度下，规模收益规律是否具备跨病种的外部有效性？
\end{enumerate}

为了回应上述挑战，本章严格依据比较统计学原则，引入两组定位截然不同、独立互补的外部验证场景：
\begin{itemize}
    \item **场景一：CDC FluSight 官方前瞻多模型集成预测（$N = 1{,}177$）**——定位为**“独立预测模型与工业级预警系统验证”**。数据源虽然同属流感监测，但其预测值是由全美 20 余个顶级建模团队（包括哈佛大学、哥伦比亚大学、洛斯阿拉莫斯国家实验室等）独立提交、官方加权集成的真实前瞻多步预报；
    \item **场景二：欧洲 11 国早期 COVID-19 跨国日度监测（$N = 404$）**——定位为**“独立病种、独立地理区域与日度高频尺度验证”**。拓展至新型冠状病毒（SARS-CoV-2）在大流行初期的跨国日度病例流。
\end{itemize}

\section{CDC FluSight 官方多模型前瞻集成预测验证}\label{sec:flusight}

\subsection{数据源架构与纯事前样本筛选规范}
美国 CDC 流感预测协作网（FluSight）自 2013 年起启动运行。在 2022--2023 流感季，FluSight 汇聚了全美 22 个独立建模团队提交的分周前瞻预测，并通过自适应非负分位数加权构建出官方集成预测（Ensemble Forecast）\citep{mathis2024collaborative}。

在以往的数据处理中，部分外部研究常回溯利用“未来实际住院数 $\ge 5$”来筛选预测样本。本文严格指出：**在决策当期，未来的实际值尚未发生，任何依赖未来真实结果的样本筛选都构成隐性数据穿越与前瞻偏误（Look-ahead Bias）**。
本文在 FluSight 实证中严格推行**纯事前信息集筛选（Strictly Ex-ante Filtering）**：
\begin{itemize}
    \item 筛选门槛仅依赖预测发布时刻已知的集成点预测值：$\widehat{C}_{s, t+1 \mid t} \ge 5$；
    \item 严格区分地点层级，排除汇总的全国序列与波多黎各海外属地，仅保留全美 50 州与华盛顿特区；
    \item 最终获得完全基于事前信息集构建的 $N = 1{,}177$ 个有效集成预测评估样本。
\end{itemize}

\subsection{集成预测实证估计结果}
对 FluSight 集成相对预测损失应用 FGLS 拟合，并执行 1,500 次按州聚类的 Bootstrap 抽样，估计结果汇报于表~\ref{tab:flusight_results}。

\begin{table}[htbp]
\centering
\small
\caption{CDC FluSight 官方多模型前瞻集成预测实证估计与敏感性对照表}
\label{tab:flusight_results}
\begin{tabularx}{\textwidth}{lcccccX}
\toprule
\textbf{样本筛选口径} & \textbf{样本量 $N$} & \textbf{截距 $\hat{a}$} & \textbf{规模项 $\hat{b}$} & \textbf{平衡点 $m_\times$} & \textbf{95\% 聚类 Bootstrap CI} & \textbf{流行病学与模型含义} \\
\midrule
\textbf{纯事前口径 (基准)} & $\mathbf{1{,}177}$ & $\mathbf{0.1294}$ & $\mathbf{7.8045}$ & $\mathbf{60.31}$ & $[5.3582, \; 10.5260]$ & **仅依据事前预测值筛选，无数据穿越** \\
回顾性事后对照口径 & $1{,}177$ & $0.1164$ & $8.8081$ & $75.67$ & $[5.9120, \; 12.1100]$ & 额外要求未来实际住院 $\ge 5$ \\
全美全辖区含属地 & $1{,}215$ & $0.1288$ & $7.7210$ & $59.95$ & $[5.2890, \; 10.4120]$ & 纳入波多黎各与境外属地后完全稳健 \\
\bottomrule
\end{tabularx}
\end{table}

实证估计显示：
\begin{enumerate}[label=(\arabic*)]
    \item **规模反比项高度显著为正**：在真实的工业级前瞻多模型集成预测下，规模系数估计值为 $\hat{b} = 7.8045$，其州聚类 95\% 置信区间为 $[5.3582, 10.5260]$，证实即使采用了复杂的机器学习与动力学集成算法，预测误差随病例规模反比衰减的基本定律依然牢固成立；
    \item **经验平衡点向右偏移至 60 例/周**：经验平衡点估计值为 $m_\times \approx 60.3$ 例/周（95\% CI $[34.5, 95.9]$），相较于第五章单基准的 22.6 例明显右移。这一位移反映了多模型集成有效吸收了部分宏观模型设定误差，使得环境方差底板相对微观抽样项的占比发生了动态调整。
\end{enumerate}

\subsection{达峰期均值偏差剧烈消除：多模型集成的终极价值}
最为惊人的实证发现在于达峰期的均值--方差分解。
在第五章中，采用单一平稳基准预测时，达峰期看似恒定的高误差中整整 88.9\% 源于宏观基准漂移。
而在 CDC FluSight 官方多模型前瞻集成系统下，重新执行精确的均值--方差相加分解（严格采用自由度 $\text{ddof} = 0$ 计算样本总体方差与均值偏差平方）：
\begin{itemize}
    \item 流行上升期：平均相对偏差平方贡献占总均方误差的 **32.8\%**；
    \item **达峰与转折期**：平均相对偏差平方贡献占总均方误差的比例从单一基准的 **88.9\% 骤降至 15.5\%**！
    \item 回落期：偏差平方占比进一步萎缩至 **0.28\%**。
\end{itemize}

这一跨模型验证给出了极其震撼的科学结论：
\begin{center}
\fbox{\parbox{0.92\textwidth}{
\textbf{第五章中高达 88.9\% 的达峰常数偏差，在多模型集成预测下被剧烈压缩至 15.5\%！\\
这在真实世界中直接证实了：跨过特征平衡点 $m_\times$ 后，继续扩大样本捕获无法克服环境平台；但转向多模型前瞻自适应集成，能够以极高的效率校准系统均值漂移，实现预测精度的阶跃式突破！}
}}
\end{center}

\section{欧洲 11 国 COVID-19 早期传播跨病种扩展验证}\label{sec:covid11}

\subsection{欧洲 11 国早期大流行日度监测数据背景}
为了检验仿射规律在不同病原体、不同国家体制以及日度高频观测下的泛化能力，本文引入 Flaxman 等人 \citep{flaxman2020estimating} 发表于《Nature》的欧洲 11 国（英国、意大利、西班牙、法国、德国、比利时、瑞士、瑞典、荷兰、挪威、奥地利）在 2020 年 2 月至 4 月大流行初期的日度发病与住院监测数据。

由于大流行早期各国经历了不同强度的封城隔离（NPIs），有效再生数随时间剧烈波动，且数据按自然日结算，构成了与全美周度流感完全正交的全新检验场景。

\subsection{纯事前增长样本筛选与参数估计}
在日度高频数据中，真实有效暴露规模 $D_t$ 由既往 7 天移动卷积和给出。为了避免数据穿越，本文同样采用纯事前筛选：保留历史规模 $D_t \ge 5$ 且前期历史增长比 $g > 0$ 的有效国家--日观测，获得分析样本量 $N = 404$。

表~\ref{tab:covid_results} 汇报了欧洲 COVID-19 实证估计及不同筛选口径下的敏感性对照。

\begin{table}[htbp]
\centering
\small
\caption{欧洲 11 国早期 COVID-19 日度监测实证估计与口径敏感性对比}
\label{tab:covid_results}
\begin{tabularx}{\textwidth}{lcccccX}
\toprule
\textbf{样本筛选口径} & \textbf{样本量 $N$} & \textbf{截距 $\hat{a}$} & \textbf{规模项 $\hat{b}$} & \textbf{日度平衡点 $m_\times$} & \textbf{95\% 国家聚类 CI} & \textbf{敏感性评估与科学启示} \\
\midrule
\textbf{纯事前有效规模 (基准)} & $\mathbf{404}$ & $\mathbf{0.2720}$ & $\mathbf{33.9223}$ & $\mathbf{124.71}$ & $[14.6021, \; 74.9873]$ & **完全事前信息集，规模项显著为正** \\
保留全增长期 ($g > 1$) & $318$ & $0.2103$ & $18.5280$ & $88.10$ & $[11.2310, \; 42.1500]$ & 斜率保持显著为正，拐点降至 88 例/天 \\
原文献截断子集 ($1 < g < 10$) & $307$ & $0.1886$ & $12.7632$ & $67.67$ & $[8.9510, \; 28.4100]$ & 排除右尾极大值后平衡点降至 67 例/天 \\
\bottomrule
\end{tabularx}
\end{table}

实证检验展现出以下核心结论：
\begin{enumerate}[label=(\arabic*)]
    \item **仿射反比递减形态在新型冠状病毒日度数据中依然坚挺**：在纯事前基准口径下，规模反比项系数为 $\hat{b} = 33.9223$，国家聚类 95\% 置信区间为 $[14.6021, 74.9873]$，显著远离零点，有力支持了规模收益递减规律跨越不同病毒病种的普遍性；
    \item **经验平衡点与时间分辨率的强依赖性**：欧洲 COVID-19 的日度经验平衡点处于 **67.7 至 124.7 例/日** 区间。与全美流感周度尺度（约 22.6 例/周）相比，日度平衡点显著更高。这一差异在动力学上极其合理：日度数据的抽样稀疏度远高于周度聚合数据，微观泊松相遇噪声与单日上报随机性更加剧烈，因此需要更大的单日病例规模才能压平抽样噪声逼近不可约环境底板。
\end{enumerate}

\section{外部验证的核心启示与经验边界界定}\label{sec:ext_summary}

两组具有高度异质性的扩展验证场景，为全文的理论框架提供了极为坚实的支撑：
\begin{center}
\fbox{\parbox{0.92\textwidth}{
\textbf{两组扩展场景共同证实：在不同的独立预测模型、跨国别行政体制与跨病种大流行监测中，均一致观察到了规模相关误差项显著为正（$b > 0$）的核心规律；\\
同时，外部检验严格划定了理论应用的方法学边界：经验平衡点 $m_\times$ 并不是一个跨病种通用的机械物理常数，其具体取值敏感地依赖于监测的时间粒度、报告流程延迟以及病原体的内在传播强度。}
}}
\end{center}
这一结论在流行病学方法学上具有重要意义：公共卫生部门在应用本理论时，应当将 $m_\times$ 视为具体监测系统在特定阶段下的**量级参照与运筹判据**，而非僵化套用单一数字。在下一章中，我们将系统探讨这一理论对国家公共卫生监测体系资源优化的政策建议。
